\section{Soybean field trial construction and analysis}\label{supp:soynam}
\newcommand{\SoyLines}{5533}
\newcommand{\SoyFamilies}{40}
\newcommand{\SoyRecords}{11066}
\newcommand{\SoyOriginalMarkers}{4240}
\newcommand{\SoyOriginalLines}{5555}
\newcommand{\SoyOriginalRecords}{12762}
\newcommand{\SoyMarkers}{4235}
\newcommand{\SoyDimension}{8526}
\newcommand{\SoyCheckPlots}{1643}
\newcommand{\SoyTestLines}{1107}
\newcommand{\SoyTrainingLines}{4426}
\newcommand{\SoyYieldMeanThirteen}{74.78}
\newcommand{\SoyYieldMeanFourteen}{55.73}
\newcommand{\SoyFits}{21}
\newcommand{\SoyMinSelected}{100}
\newcommand{\SoyMaxSelected}{181}
\newcommand{\SoyMaxKKT}{9.09e-09}
\newcommand{\SoyMaxNuisance}{4.73e-11}
\newcommand{\SoyReferenceMinMinutes}{9.5}
\newcommand{\SoyReferenceMaxMinutes}{11.9}
\newcommand{\SoyMedianThirteen}{3.96}
\newcommand{\SoyMedianThirteenInterval}{$[3.01,\ 4.91]$}
\newcommand{\SoyMedianFourteen}{5.16}
\newcommand{\SoyMedianFourteenInterval}{$[4.19,\ 6.13]$}
\newcommand{\SoyUpperYearDifference}{1.61}
\newcommand{\SoyUpperYearDifferenceInterval}{$[0.14,\ 3.07]$}
\newcommand{\SoyLowerYearDifference}{-0.19}
\newcommand{\SoyLowerYearDifferenceInterval}{$[-1.86,\ 1.47]$}
\newcommand{\SoyMedianYearDifference}{1.20}
\newcommand{\SoyMedianYearDifferenceInterval}{$[-0.05,\ 2.45]$}
\newcommand{\SoyBandwidthChange}{0.010}
\newcommand{\SoyCoefficientMinimum}{3.85}
\newcommand{\SoyCoefficientMaximum}{5.45}
\newcommand{\SoyYearContrastMinimum}{1.50}
\newcommand{\SoyYearContrastMaximum}{2.02}
\newcommand{\SoySpatialYearContrast}{2.02}
\newcommand{\SoySpatialYearContrastInterval}{$[0.40,\ 3.63]$}
\newcommand{\SoyGenomicLowerGain}{1.5}
\newcommand{\SoyGenomicLowerGainInterval}{$[-0.4,\ 3.4]$}
\newcommand{\SoyGenomicLowerLossDifference}{0.041}
\newcommand{\SoyGenomicMedianGain}{1.6}
\newcommand{\SoyGenomicMedianGainInterval}{$[-0.2,\ 3.4]$}
\newcommand{\SoyGenomicMedianLossDifference}{0.052}
\newcommand{\SoyGenomicUpperGain}{2.3}
\newcommand{\SoyGenomicUpperGainInterval}{$[0.6,\ 3.9]$}
\newcommand{\SoyGenomicUpperLossDifference}{0.059}


\subsection{Released data and plot identifiers}

The analysis uses \texttt{NAM} version 1.8.0, object \texttt{met},
and \texttt{SoyNAM} version 1.6.2, object \texttt{soyin}. The first
release supplies the recorded canopy values, yields and genotype
matrix. The second supplies family labels, field check plots and the
yield-unit dictionary. The two releases agree on yield to within
0.05 bu/ac at all 11088 matched 2013--2014 plot coordinates. Maturity
records agree exactly. Plot coordinates identify 277 records whose
strain-name prefixes differ between releases; the analysis retains
the \texttt{NAM} identifiers that match its genotype rows.

The released panel contains \SoyOriginalLines{} lines and
\SoyOriginalRecords{} records across three years. We use the two
years in which the complete family panel was planted
\cite{xavier2017canopy}. Requiring a complete pair retains
\SoyLines{} lines and \SoyRecords{} observations. The 2015 trial
contains selected families under a different sampling design.
The five duplicate marker columns are removed before analysis;
all remaining \SoyMarkers{} released genotype columns are observed.
The deposited matrix supplies the processed genotype calls.

\begin{table}[t]
\centering
\small
\caption{Observed distributions in the complete two-year panel. Each year contains one observation for each of the same 5533 lines.}
\label{tab:soy-descriptive}
\begin{tabular}{rlrrrrr}
\toprule
Year & Variable & Mean & SD & Quartile 1 & Median & Quartile 3 \\
\midrule
2013 & Yield (bu/ac) & 74.78 & 10.46 & 68.51 & 75.43 & 81.98 \\
2014 & Yield (bu/ac) & 55.73 & 11.57 & 47.80 & 55.60 & 63.38 \\
2013 & Canopy coverage (\%) & 66.63 & 3.80 & 64.25 & 66.95 & 69.26 \\
2014 & Canopy coverage (\%) & 56.71 & 4.91 & 53.82 & 57.16 & 60.08 \\
2013 & Maturity (days) & 128.87 & 5.40 & 126.00 & 130.00 & 132.00 \\
2014 & Maturity (days) & 126.39 & 6.13 & 123.00 & 128.00 & 130.00 \\
\bottomrule
\end{tabular}
\end{table}


Yield is divided by ten for calculation and restored to bu/ac in
every reported coefficient, interval and prediction loss. The working
penalty $\Lambda=1$ is defined on this response scale, equivalently
$\Lambda=0.1$ when the complete check-loss criterion is expressed in
bu/ac. The two
canopy columns are $ACC_{ij}1\{Year_{ij}=2013\}/10$ and
$ACC_{ij}1\{Year_{ij}=2014\}/10$. Each SNP dosage similarly has two
year-specific columns. Family adjustment uses 39 contrasts; the eight
block indicators include the year intercepts. Maturity uses natural
cubic splines with knots at 120 and 130 days and boundary knots at
100 and 160 days, interacted with year. The reference analysis retains
all recorded maturity values, including the two values above 180 days.

\subsection{Check plot calibration and sample roles}

The spatial calibration includes \SoyCheckPlots{} check plots and
90 check-strain--year combinations. It uses a Gaussian additive model
with strain--year and block effects, and a tensor product cubic
regression spline in row and column within each block. Each marginal
basis has dimension five. REML chooses the smoothing parameters,
with shrinkage permitted. Only the fitted spatial terms supply the
analysed lines' check-surface covariate. Physical plot coordinates
verify that these check plots and the analysed line plots are disjoint.
The inference conditions on the supplied calibration surface.

Seed 20261031 orders lines randomly within each family. Alternating
positions define the two complementary inference samples. Three
inner groups within each selection sample choose the outcome-search
size and the profile-adjustment penalties. Training samples alone
supply the scales used by their numerical subproblems. The complete
genomic block remains eligible for selection at each stage.

A separate cyclic allocation assigns approximately one fifth of
each family's lines to the prediction test set. Both observations
of a line have the same role. This yields \SoyTrainingLines{} training
and \SoyTestLines{} test lines. Internal prediction validation uses
three complete-line groups from the training set. The two inference
samples and the prediction training sample serve different statistical
calculations; test outcomes never enter prediction selection or
calibration.

\subsection{Certified convex fits for paired observations}

The outcome search uses a weighted orthogonal basis. For a selected
design $U$ and line-specific intercepts $\gamma_i$, its convex problem,
before division by the number of lines, is
\[
 \min_{a,\gamma}\left\{
 \frac12\lVert a\rVert^2+\frac\lambda2\sum_i\gamma_i^2
 +\frac12\sum_{i,j}\rho_\tau(Y_{ij}-U_{ij}^Ta-\gamma_i)
 \right\}.
\]
The coefficient ridge equals the registered prediction ridge under
the weighted change of coordinates. Let $W=I/2$ and let $C$ contain
the within-line row means of $U$ as columns. In a two-block ADMM
update with augmentation $\eta$, elimination of the line intercepts
gives the coefficient matrix
\[
 I+\eta U^TWU-\frac{\eta^2}{\lambda+\eta}CC^T.
\]
Its dimension is the selected coefficient dimension, and its
factorisation is shared by all iterations of that selected fit.
The check-loss proximal map is exact. The implementation fixes
$\eta=1$ and relaxation 1.6. Equality-constrained refinement on the
residual active face supplies a candidate solution when its linear
system is identified.

For dual scores $s_{ij}\in[(\tau-1)/2,\tau/2]$, the dual value is
\[
 D(s)=Y^Ts-\frac12\lVert U^Ts\rVert^2
      -\frac1{2\lambda}\sum_i(s_{i1}+s_{i2})^2.
\]
Every returned candidate satisfies the original primal--dual gap
criterion after division by the number of lines. Independent
Clarabel calculations on actual observations check the objective
and fitted values across all three quantiles and penalties
$\lambda=0.5,1,2$, with additional selected dimensions up to 419.

The final smoothed fits use weighted QR coordinates. Writing
$W_0^{1/2}X_AP=QR$ with $W_{0,jj}=1/(2n_f)$ gives
$U=W_0^{-1/2}Q$ and $a=RP^T\beta_A$. Thus $Ua=X_A\beta_A$ and
the prediction ridge becomes $a^Ta$. Newton iteration retains the
complete profile curvature. Additional precision iterations enforce
the original-coordinate gradient criterion after transforming back.
The maximum original-coordinate KKT residual among the reported
systems is \texttt{\SoyMaxKKT{}}.

\subsection{Joint coordinate and contrast intervals}

Collect the six year-specific quantile estimates in $\widehat\theta$.
Line identifiers align their saved influence vectors $V_i$ across
the three quantiles. Within each estimation group $I_f$, the
contributions are centred before constructing
\[
 \widehat\Sigma=\frac1{n^2}\sum_f\sum_{i\in I_f}
 (V_i-\bar V_f)(V_i-\bar V_f)^T.
\]
The diagonal reproduces the six reported standard errors. A contrast
$c$ has estimate $c^T\widehat\theta$ and standard error
$(c^T\widehat\Sigma c)^{1/2}$. The six coordinates, two within-year
upper-minus-lower contrasts and three between-year contrasts share
critical value $z_{1-0.05/(2\times11)}$.

\begin{table}[t]
\centering
\small
\caption{Canopy profile coefficients in bu/ac per ten percentage points of coverage. Simultaneous intervals allocate 5\% across six coordinates and five contrasts. The reference penalty is $\Lambda=1$.}
\label{tab:soy-reference}
\begin{tabular}{rrrrr}
\toprule
Year & $\tau$ & Estimate & Standard error & Simultaneous interval \\
\midrule
2013 & 0.25 & 4.542 & 0.434 & $[3.310,\ 5.774]$ \\
2013 & 0.50 & 3.959 & 0.334 & $[3.011,\ 4.908]$ \\
2013 & 0.75 & 3.845 & 0.373 & $[2.787,\ 4.904]$ \\
2014 & 0.25 & 4.348 & 0.430 & $[3.127,\ 5.569]$ \\
2014 & 0.50 & 5.157 & 0.342 & $[4.185,\ 6.128]$ \\
2014 & 0.75 & 5.452 & 0.381 & $[4.370,\ 6.534]$ \\
\bottomrule
\end{tabular}
\end{table}

\begin{table}[t]
\centering
\small
\caption{Prespecified canopy coefficient differences in bu/ac per ten percentage points. These intervals use the same eleven-parameter family as the coordinate intervals.}
\label{tab:soy-contrasts}
\begin{tabular}{lrrr}
\toprule
Contrast & Estimate & Standard error & Simultaneous interval \\
\midrule
Upper minus lower in 2013 & -0.697 & 0.467 & $[-2.022,\ 0.629]$ \\
Upper minus lower in 2014 & 1.104 & 0.482 & $[-0.264,\ 2.471]$ \\
2014 minus 2013 at $\tau=0.25$ & -0.194 & 0.588 & $[-1.863,\ 1.475]$ \\
2014 minus 2013 at $\tau=0.50$ & 1.198 & 0.441 & $[-0.055,\ 2.450]$ \\
2014 minus 2013 at $\tau=0.75$ & 1.606 & 0.517 & $[0.138,\ 3.075]$ \\
\bottomrule
\end{tabular}
\end{table}


\subsection{Sensitivity calculations}

All seven settings use the same three quantiles. Bandwidth
multipliers 0.75 and 1.25 retain the reference target. The settings
$\Lambda=0.5$ and $2$ define their respective profile parameters.
The maturity setting excludes the complete pair for each of the two
lines with maturity above 180 days. The spatial setting adds row,
column, squared row, squared column and their product separately
within each block. Its scales are fixed functions of plot position.

\begin{table}[t]
\centering
\footnotesize
\caption{Complete registered sensitivity calculations. Entries are estimates and simultaneous intervals in bu/ac per ten percentage points. Each setting uses the eleven-parameter critical value; the working penalty belongs to its profile target.}
\label{tab:soy-sensitivity}
\begin{tabular}{lrrr}
\toprule
Setting & $\tau$ & 2013 & 2014 \\
\midrule
Reference & 0.25 & 4.542 $[3.310,\ 5.774]$ & 4.348 $[3.127,\ 5.569]$ \\
Reference & 0.50 & 3.959 $[3.011,\ 4.908]$ & 5.157 $[4.185,\ 6.128]$ \\
Reference & 0.75 & 3.845 $[2.787,\ 4.904]$ & 5.452 $[4.370,\ 6.534]$ \\
$0.75h$ & 0.25 & 4.543 $[3.262,\ 5.824]$ & 4.348 $[3.183,\ 5.512]$ \\
$0.75h$ & 0.50 & 3.960 $[3.024,\ 4.897]$ & 5.147 $[4.149,\ 6.146]$ \\
$0.75h$ & 0.75 & 3.846 $[2.765,\ 4.927]$ & 5.457 $[4.371,\ 6.543]$ \\
$1.25h$ & 0.25 & 4.536 $[3.313,\ 5.760]$ & 4.350 $[3.134,\ 5.567]$ \\
$1.25h$ & 0.50 & 3.962 $[3.004,\ 4.920]$ & 5.163 $[4.226,\ 6.101]$ \\
$1.25h$ & 0.75 & 3.847 $[2.782,\ 4.912]$ & 5.442 $[4.375,\ 6.508]$ \\
$\Lambda=0.5$ & 0.25 & 4.504 $[3.450,\ 5.558]$ & 4.235 $[3.104,\ 5.366]$ \\
$\Lambda=0.5$ & 0.50 & 3.959 $[2.828,\ 5.091]$ & 5.108 $[4.048,\ 6.167]$ \\
$\Lambda=0.5$ & 0.75 & 3.892 $[2.865,\ 4.919]$ & 5.391 $[4.387,\ 6.394]$ \\
$\Lambda=2$ & 0.25 & 4.573 $[3.233,\ 5.912]$ & 4.394 $[2.767,\ 6.020]$ \\
$\Lambda=2$ & 0.50 & 4.073 $[3.147,\ 4.999]$ & 4.976 $[4.012,\ 5.939]$ \\
$\Lambda=2$ & 0.75 & 3.926 $[2.875,\ 4.976]$ & 5.529 $[4.424,\ 6.634]$ \\
Maturity restriction & 0.25 & 4.636 $[3.462,\ 5.809]$ & 4.149 $[2.999,\ 5.299]$ \\
Maturity restriction & 0.50 & 4.140 $[3.059,\ 5.222]$ & 4.730 $[3.675,\ 5.785]$ \\
Maturity restriction & 0.75 & 3.682 $[2.486,\ 4.878]$ & 5.424 $[4.506,\ 6.343]$ \\
Spatial polynomial & 0.25 & 4.794 $[3.646,\ 5.942]$ & 5.261 $[4.085,\ 6.438]$ \\
Spatial polynomial & 0.50 & 4.222 $[3.150,\ 5.293]$ & 5.482 $[4.408,\ 6.555]$ \\
Spatial polynomial & 0.75 & 3.844 $[2.704,\ 4.983]$ & 5.860 $[4.664,\ 7.056]$ \\
\bottomrule
\end{tabular}
\end{table}

\begin{table}[t]
\centering
\small
\caption{The 2014 minus 2013 canopy coefficient contrast at the upper quartile across all registered settings. Units are bu/ac per ten percentage points of coverage; intervals use the eleven-parameter critical value within each setting.}
\label{tab:soy-contrast-sensitivity}
\begin{tabular}{lrrr}
\toprule
Setting & Estimate & Standard error & Simultaneous interval \\
\midrule
Reference & 1.606 & 0.517 & $[0.138,\ 3.075]$ \\
$0.75h$ & 1.611 & 0.525 & $[0.121,\ 3.101]$ \\
$1.25h$ & 1.595 & 0.515 & $[0.133,\ 3.057]$ \\
$\Lambda=0.5$ & 1.499 & 0.465 & $[0.178,\ 2.819]$ \\
$\Lambda=2$ & 1.603 & 0.520 & $[0.126,\ 3.080]$ \\
Maturity restriction & 1.742 & 0.500 & $[0.322,\ 3.162]$ \\
Spatial polynomial & 2.016 & 0.568 & $[0.404,\ 3.628]$ \\
\bottomrule
\end{tabular}
\end{table}


\subsection{Prediction methods and uncertainty}

Prediction compares the joint profile with canopy, the same procedure
without canopy, and pooled quantile Lasso with canopy. Family and
field adjustments are common to all methods. Maturity is excluded.
Each joint predictor uses the unsmoothed convex outcome-search path
with $\Lambda=1$, selecting the model size by complete-line validation
pinball loss. The candidate sizes range from zero to
$\min_f\lfloor n_{train,f}/\log(n_{train,f})\rfloor$ additional
genomic variables beyond the retained adjustment terms. The final predictor uses the fixed-effect component for
a line absent from training.

The pooled method minimises the complete check loss with a Lasso
penalty on standardised genomic variables; the prespecified
adjustment variables are unpenalised. Each training sample uses
40 geometric penalty ratios from 1 to 0.01. The entry value is the
largest absolute standardised genomic score from a dual solution of
the unpenalised adjustment model, divided by the number of observations.
HiGHS solves the dual linear programme, retaining its basis along
the path. Constraints are added whenever a complete genomic scan
finds a violation. Every returned point has normalised primal--dual
gap at most $10^{-7}$ and dual-constraint error at most $10^{-8}$.
An independent \texttt{quantreg} calculation checks the objective,
including the zero-penalty limit. Both fitting and evaluation use
the unsmoothed check loss.

For each method and year, the empirical $\tau$ quantile of internal
validation residuals supplies an intercept calibration. We use the
inverse empirical distribution, R quantile type 1. The fitted model
and calibration are saved before test losses are calculated.
Raw and calibrated results are both retained. The calibrated
comparison is the registered prediction assessment.

Let $A_i$ and $B_i$ be the two-year average check losses for the
canopy joint predictor and a comparator on test line $i$. The
reported improvement is $\bar B-\bar A$, with standard error
$s_{B-A}/\sqrt{n_{test}}$. Relative improvement is
$100(1-\bar A/\bar B)$; its delta-method standard error uses the
paired line contributions
\[
 100\left\{-\frac{A_i-\bar A}{\bar B}
 +\frac{\bar A(B_i-\bar B)}{\bar B^2}\right\}.
\]
The three quantiles and two method comparisons use critical value
$z_{1-0.05/(2\times6)}$. These intervals describe held-out-line
variation conditional on the training models and field calibration.

\begin{table}[t]
\centering
\small
\caption{Mean pinball loss on the same held-out lines, in bu/ac. Each line receives equal weight across its two observations. Every predictor uses calibration estimated from training validation residuals.}
\label{tab:soy-prediction}
\begin{tabular}{rrr}
\toprule
$\tau$ & Joint with canopy & Joint without canopy \\
\midrule
0.25 & 2.670 & 2.711 \\
0.50 & 3.240 & 3.292 \\
0.75 & 2.545 & 2.604 \\
\bottomrule
\end{tabular}
\end{table}

\begin{table}[t]
\centering
\footnotesize
\caption{Raw and training-calibrated predictions on held-out lines. Coverage is the fraction of observed yields at or below the predicted quantile. The nominal fractions are the corresponding values of $\tau$.}
\label{tab:soy-prediction-details}
\begin{tabular}{lrrrrr}
\toprule
Method & $\tau$ & Raw loss & Calibrated loss & Raw coverage & Calibrated coverage \\
\midrule
Joint profile with canopy & 0.25 & 2.704 & 2.670 & 0.301 & 0.269 \\
Joint profile with canopy & 0.50 & 3.240 & 3.240 & 0.501 & 0.503 \\
Joint profile with canopy & 0.75 & 2.550 & 2.545 & 0.714 & 0.773 \\
Joint profile without canopy & 0.25 & 2.740 & 2.711 & 0.298 & 0.262 \\
Joint profile without canopy & 0.50 & 3.291 & 3.292 & 0.514 & 0.517 \\
Joint profile without canopy & 0.75 & 2.606 & 2.604 & 0.728 & 0.768 \\
\bottomrule
\end{tabular}
\end{table}


\bibliography{references}
\end{document}
